\documentclass[10pt, aspectratio=169]{beamer}
% Preámbulo modular para presentaciones Beamer (Modelación Estadística)
% Diseño y paleta institucional del Tecnológico de Monterrey

\documentclass[aspectratio=169,xcolor={dvipsnames,table}]{beamer}

\usepackage[utf8]{inputenc}
\usepackage[spanish,mexico]{babel}
\usepackage{amsmath,amssymb,amsfonts,mathtools}
\usepackage{graphicx}
\usepackage{listings}
\usepackage{booktabs}
\usepackage{tikz}
\usetikzlibrary{matrix,arrows,decorations.pathmorphing,shapes.geometric,positioning}

% Paleta Institucional Tecnológico de Monterrey
\definecolor{TecAzul}{HTML}{0039A6}       % Azul TEC: color primario institucional
\definecolor{TecAzulOscuro}{HTML}{1A2E51} % Azul marino oscuro
\definecolor{TecRojo}{HTML}{EC2661}       % Rojo de acento (paleta miTec)
\definecolor{TecGris}{HTML}{646464}       % Gris neutro para comentarios y subtítulos
\definecolor{TecFondoBloque}{HTML}{F4F6F9}% Fondo suave para bloques

% Configuración de Tema Beamer
\usetheme{Boadilla}
\usecolortheme[named=TecAzulOscuro]{structure}
\setbeamercolor{alerted text}{fg=TecRojo}
\setbeamercolor{palette primary}{bg=TecAzulOscuro,fg=white}
\setbeamercolor{palette secondary}{bg=TecAzul,fg=white}
\setbeamercolor{palette tertiary}{bg=TecAzulOscuro!80!black,fg=white}
\setbeamercolor{title}{fg=TecAzulOscuro,bg=TecFondoBloque}
\setbeamercolor{frametitle}{fg=TecAzulOscuro,bg=TecFondoBloque}

% Bloques personalizados elegantes
\setbeamercolor{block title}{bg=TecAzulOscuro,fg=white}
\setbeamercolor{block body}{bg=TecFondoBloque,fg=black}
\setbeamercolor{block title alerted}{bg=TecRojo,fg=white}
\setbeamercolor{block body alerted}{bg=TecRojo!8,fg=black}
\setbeamercolor{block title example}{bg=TecAzul,fg=white}
\setbeamercolor{block body example}{bg=TecAzul!8,fg=black}

% Redondear bordes y añadir sombra sutil
\setbeamertemplate{blocks}[rounded][shadow=true]
\setbeamertemplate{navigation symbols}{}

% Pie de página limpio con número de diapositiva
\setbeamertemplate{footline}{
  \leavevmode%
  \hbox{%
  \begin{beamercolorbox}[wd=.7\paperwidth,ht=2.25ex,dp=1ex,left]{author in head/foot}%
    \hspace*{2ex}\insertshortauthor~~(\insertshortinstitute) \hfill \insertshorttitle
  \end{beamercolorbox}%
  \begin{beamercolorbox}[wd=.3\paperwidth,ht=2.25ex,dp=1ex,right]{date in head/foot}%
    \insertshortdate{}\hspace*{2em}
    \insertframenumber{} / \inserttotalframenumber\hspace*{2ex}
  \end{beamercolorbox}}%
  \vskip0pt%
}

% Configuración de Listings de Python para visualización pedagógica de código
\lstset{
    language=Python,
    basicstyle=\ttfamily\footnotesize,
    keywordstyle=\color{TecAzulOscuro}\bfseries,
    stringstyle=\color{TecRojo},
    commentstyle=\color{TecGris}\itshape,
    showstringspaces=false,
    breaklines=true,
    frame=lines,
    rulecolor=\color{TecAzulOscuro!30},
    backgroundcolor=\color{TecFondoBloque},
    aboveskip=0.5em,
    belowskip=0.5em,
    literate={á}{{\'a}}1 {é}{{\'e}}1 {í}{{\'i}}1 {ó}{{\'o}}1 {ú}{{\'u}}1
             {Á}{{\'A}}1 {É}{{\'E}}1 {Í}{{\'I}}1 {Ó}{{\'O}}1 {Ú}{{\'U}}1
             {ñ}{{\~n}}1 {Ñ}{{\~N}}1 {¿}{{?`}}1 {¡}{{!`}}1
}

% Metadatos institucionales por defecto
\title[Modelación Estadística]{Modelación Estadística para Ciencia de Datos e Ingeniería}
\author[J. Castillo Colmenares]{Juliho David Castillo Colmenares}
\institute[Tec de Monterrey]{Tecnológico de Monterrey \\ Escuela de Ingeniería y Ciencias}
\date{\today}

% Comandos y macros estadísticas compartidas para las presentaciones Beamer (Metropolis)

% Conjuntos numéricos básicos
\newcommand{\R}{\mathbb{R}}
\newcommand{\N}{\mathbb{N}}
\newcommand{\Q}{\mathbb{Q}}
\newcommand{\Z}{\mathbb{Z}}
\newcommand{\set}[1]{\left\{ #1 \right\}}
\newcommand{\abs}[1]{\left|#1\right|}
\newcommand{\norm}[1]{\left\| #1 \right\|}

% Comandos estadísticos y probabilísticos del curso
\newcommand{\Var}{\operatorname{Var}}
\newcommand{\cov}{\operatorname{Cov}}
\newcommand{\corr}{\rho}
\newcommand{\comb}[2]{\begin{pmatrix} #1 \\ #2 \end{pmatrix}}
\newcommand{\s}{\sigma}
\newcommand{\particion}{\mathfrak{P}}
\newcommand{\card}[1]{\#\left( #1 \right)}
\newcommand{\seq}[3]{\set{#1}_{#2}^{#3}}
\newcommand{\E}{\mathbb{E}}
\newcommand{\Prob}{\mathbb{P}}

% Entornos pedagógicos y cajas en español académico natural
\newenvironment{discusion}{%
  \begin{alertblock}{Pregunta de Discusión}%
}{%
  \end{alertblock}%
}

\newenvironment{reflexion}{%
  \begin{alertblock}{Reflexión para el Aula}%
}{%
  \end{alertblock}%
}

\newenvironment{paradoja}{%
  \begin{alertblock}{Paradoja de Exploración}%
}{%
  \end{alertblock}%
}

\newenvironment{motivacion}{%
  \begin{exampleblock}{Motivación: Ciencia de Datos en la Práctica}%
}{%
  \end{exampleblock}%
}

\newenvironment{retoclase}{%
  \begin{block}{Reto Rápido en Equipo}%
}{%
  \end{block}%
}


\title[Chi-Cuadrada y Varianza Muestral]{Sección 05.03: Distribución Chi-Cuadrada y Varianza Muestral}
\subtitle{Teorema de Fisher, Independencia de $\bar{X}$ y $S^2$, e Intervalos para $\sigma^2$}
\author[J. Castillo Colmenares]{Juliho Castillo Colmenares}
\institute{Tecnológico de Monterrey}
\date{\vspace{-1.2cm}}

\begin{document}

% Slide 1: Portada
\begin{frame}[plain]
  \titlepage
\end{frame}

% Slide 2: Hoja de Ruta
\begin{frame}{Hoja de Ruta --- Capítulo 05: Distribuciones de Muestreo}
  \footnotesize
  ¿Dónde nos encontramos en el desarrollo modular del capítulo? \pause
  \begin{itemize}\itemsep=0.08em
    \item \textbf{05.01} Muestreo Aleatorio Simple, Media y Varianza Muestral Insesgada \pause
    \item \textbf{05.02} Teorema del Límite Central Asintótico \pause
    \item \textbf{05.03} Distribución Chi-Cuadrada y Varianza Muestral \emph{(Hoy)} \pause
    \item \textbf{05.04} Distribución $t$ de Student y Muestras Pequeñas \pause
    \item \textbf{05.05} Distribución $F$ de Fisher-Snedecor
  \end{itemize}
  \vspace{-0.05cm}
  \begin{block}{Objetivo de la Sesión}
    Formalizar la distribución de la varianza muestral $S^2$ para poblaciones normales a través del Teorema de Fisher, demostrar la sorprendente independencia entre $\bar{X}$ y $S^2$, y construir intervalos de confianza exactos para $\sigma^2$.
  \end{block}
\end{frame}

% Slide 3: Motivación
\begin{frame}{De $\sigma$ Conocida a $\sigma$ Desconocida}
  \footnotesize
  \begin{motivacion}
    El TLC (Sección 05.02) aproxima $\bar{X}$ con una Normal cuando $\sigma$ es \emph{conocida}. En la práctica casi siempre debemos estimar $\sigma$ con $S$. Para eso necesitamos entender la distribución exacta de $S^2$ --- y aquí es donde entra la chi-cuadrada.
  \end{motivacion}

  \vspace{0.15cm}
  \pause
  \begin{itemize}\itemsep=0.1cm
    \item \textbf{Origen:} la chi-cuadrada surge al sumar cuadrados de normales estándar independientes. \pause
    \item \textbf{Conexión con 04.07:} ya vimos que $\chi^2_\nu = \text{Gamma}(\nu/2, 2)$. \pause
    \item \textbf{Pregunta central:} ¿cuál es la distribución exacta de $(n-1)S^2/\sigma^2$?
  \end{itemize}
\end{frame}

% Slide 4: Definición y Propiedades
\begin{frame}{Definición y Propiedades de la Distribución Chi-Cuadrada}
  \footnotesize
  Sean $Z_1,\ldots,Z_\nu$ independientes con $Z_i\sim N(0,1)$. \pause

  \vspace{0.1cm}
  \begin{block}{Definición y Densidad}
    $$\chi^2_\nu = \sum_{i=1}^\nu Z_i^2, \qquad f(x) = \dfrac{1}{2^{\nu/2}\Gamma(\nu/2)}x^{\nu/2-1}e^{-x/2}, \quad x>0$$
  \end{block}

  \vspace{0.1cm}
  \pause
  \begin{alertblock}{Propiedades}
    \begin{itemize}\itemsep=0.05cm
      \item $E(\chi^2_\nu)=\nu$, $\Var(\chi^2_\nu)=2\nu$; caso particular $\chi^2_\nu=\text{Gamma}(\nu/2,2)$.
      \item Reproductividad: $\chi^2_{\nu_1}+\chi^2_{\nu_2}\sim\chi^2_{\nu_1+\nu_2}$ (independientes).
    \end{itemize}
  \end{alertblock}
\end{frame}

% Slide 5: Teorema de Fisher
\begin{frame}{El Teorema de Fisher}
  \footnotesize
  Para una muestra normal iid $X_1,\ldots,X_n \sim N(\mu,\sigma^2)$: \pause

  \vspace{0.1cm}
  \begin{block}{Enunciado}
    \begin{enumerate}\itemsep=0.05cm
      \item $\bar{X}$ y $S^2$ son \textbf{independientes}.
      \item $\dfrac{(n-1)S^2}{\sigma^2} \sim \chi^2_{n-1}$.
    \end{enumerate}
  \end{block}

  \vspace{0.1cm}
  \pause
  \begin{alertblock}{¿Por Qué es Sorprendente?}
    Aunque $S^2$ se calcula a partir de $(X_i-\bar{X})$, que depende de $\bar{X}$, el resultado es independiente de $\bar{X}$. La demostración usa una transformación ortogonal: $\bar{X}$ captura una dirección, y las $n-1$ direcciones restantes (ortogonales a $\bar X$) capturan exactamente $(n-1)S^2$ --- un caso particular del \emph{Teorema de Cochran}.
  \end{alertblock}
\end{frame}

% Slide 6: Aplicación: Intervalo de Confianza para σ²
\begin{frame}{Intervalo de Confianza Exacto para $\sigma^2$}
  \footnotesize
  El Teorema de Fisher permite construir un intervalo exacto (no aproximado) para $\sigma^2$: \pause

  \vspace{0.1cm}
  \begin{block}{Fórmula del Intervalo}
    $$\left[\dfrac{(n-1)S^2}{\chi^2_{n-1,1-\alpha/2}},\; \dfrac{(n-1)S^2}{\chi^2_{n-1,\alpha/2}}\right]$$
    A diferencia de los intervalos basados en la Normal, este intervalo es \textbf{asimétrico} alrededor de $S^2$, porque la chi-cuadrada no es simétrica.
  \end{block}

  \vspace{0.1cm}
  \pause
  \begin{alertblock}{Contraste de Hipótesis para $\sigma^2$}
    El mismo estadístico $(n-1)S^2/\sigma_0^2$ permite probar $H_0:\sigma^2=\sigma_0^2$ comparándolo contra los cuantiles de $\chi^2_{n-1}$.
  \end{alertblock}
\end{frame}

% Slide 7: Aplicaciones
\begin{frame}{Aplicaciones en Inferencia y Control de Calidad}
  \footnotesize
  La distribución chi-cuadrada es central en la inferencia sobre varianzas: \pause

  \vspace{0.15cm}
  \begin{itemize}\itemsep=0.12cm
    \item \textbf{Control de calidad:} pruebas sobre la variabilidad de un proceso de manufactura ($\sigma^2$ objetivo). \pause
    \item \textbf{Bondad de ajuste:} la prueba $\chi^2$ compara frecuencias observadas vs. esperadas. \pause
    \item \textbf{Preview 05.04:} la distribución $t$ de Student se construye combinando una Normal con una chi-cuadrada independiente. \pause
    \item \textbf{Preview 05.05:} la distribución $F$ compara dos varianzas mediante el cociente de dos chi-cuadradas.
  \end{itemize}
\end{frame}

% Slide 8: Lab Python (Bloque 1)
\begin{frame}[fragile]{Laboratorio en Python: Propiedades de la Chi-Cuadrada}
  \scriptsize
  Verificación de media/varianza, conexión con la Gamma, y reproductividad (\texttt{05.03\_chi\_squared\_distribution.py}):
  {\lstset{basicstyle=\fontsize{4pt}{4.8pt}\selectfont\ttfamily,aboveskip=0.1em,belowskip=0.1em}
  \lstinputlisting[firstline=18, lastline=40]{../../code/05_distribuciones_muestreo/05.03_chi_squared_distribution.py}}
\end{frame}

% Slide 9: Lab Python (Bloque 2)
\begin{frame}[fragile]{Laboratorio en Python: Teorema de Fisher}
  \scriptsize
  Verificación Monte Carlo de la independencia de $\bar X$ y $S^2$, y de $(n-1)S^2/\sigma^2\sim\chi^2_{n-1}$:
  {\lstset{basicstyle=\fontsize{4pt}{4.8pt}\selectfont\ttfamily,aboveskip=0.1em,belowskip=0.1em}
  \lstinputlisting[firstline=43, lastline=70]{../../code/05_distribuciones_muestreo/05.03_chi_squared_distribution.py}}
\end{frame}

% Slide 10: Lab Python (Bloque 3)
\begin{frame}[fragile]{Laboratorio en Python: Cobertura del Intervalo para $\sigma^2$}
  \scriptsize
  Verificación empírica de la cobertura del intervalo de confianza chi-cuadrada para $\sigma^2$:
  {\lstset{basicstyle=\fontsize{4pt}{4.8pt}\selectfont\ttfamily,aboveskip=0.1em,belowskip=0.1em}
  \lstinputlisting[firstline=73, lastline=100]{../../code/05_distribuciones_muestreo/05.03_chi_squared_distribution.py}}
\end{frame}

% Slide 11: Salida Terminal
\begin{frame}[fragile]{Laboratorio en Python: Salida en Terminal}
  \scriptsize
  Salida estándar generada al ejecutar el script del laboratorio en Python:
  \vspace{0.02cm}
  \begin{block}{Salida Estándar en Consola}
    \fontsize{4pt}{5pt}\selectfont
    \begin{verbatim}
=== Block 1: Chi-Squared Properties ===
chi2_15: mean=15.00, variance=30.00 (matches Gamma(7.5, 2))
Additivity chi2_5+chi2_8 vs chi2_13: KS D=0.0013, p=0.908

=== Block 2: Fisher's Theorem ===
Correlation(Xbar, S^2) = 0.0038 (independence confirmed)
(n-1)*S^2/sigma^2 vs chi2_9: KS D=0.0018, p=0.517
Worked example: statistic=16.20, critical chi2_9,0.95=16.92 -> Reject H0? No

=== Block 3: Confidence Interval Coverage ===
chi2_19,0.025=8.91, chi2_19,0.975=32.85
Empirical coverage: 0.9487 (nominal: 0.9500)
95% CI for sigma^2 (n=20, S^2=45): [26.03, 96.00]
    \end{verbatim}
  \end{block}
\end{frame}

% Slide 12A: Ejercicio Nivel 1 (Enunciado)
\begin{frame}{Ejercicio en Clase --- Nivel 1: Fundamental (1/2)}
  \footnotesize
  \begin{block}{Problema 5.3.1 --- Media y Varianza de $\chi^2_{15}$ (Enunciado)}
    Sea $X \sim \chi^2_{15}$. Calcule la media $E(X)$ y la varianza $\Var(X)$.
  \end{block}

  \vspace{0.15cm}
  \pause
  \begin{alertblock}{Planteamiento}
    \begin{itemize}\itemsep=0.04cm
      \item Use directamente $E(\chi^2_\nu)=\nu$ y $\Var(\chi^2_\nu)=2\nu$.
    \end{itemize}
  \end{alertblock}
\end{frame}

% Slide 12B: Ejercicio Nivel 1 (Resolución)
\begin{frame}{Ejercicio en Clase --- Nivel 1: Fundamental (2/2)}
  \scriptsize
  Sustituimos $\nu=15$ directamente: \pause

  \vspace{0.05cm}
  \begin{block}{Cálculo}
    $$E(X) = 15, \qquad \Var(X) = 2\cdot 15 = 30.$$
  \end{block}

  \vspace{0.05cm}
  \pause
  \begin{alertblock}{Interpretación}
    \scriptsize
    A diferencia de la Normal, la chi-cuadrada tiene media y varianza que crecen juntas con $\nu$: a más grados de libertad, mayor dispersión absoluta, aunque relativamente la distribución se vuelve más simétrica (se aproxima a la Normal por el TLC).
  \end{alertblock}
\end{frame}

% Slide 13A: Ejercicio Nivel 2 (Enunciado)
\begin{frame}{Ejercicio en Clase --- Nivel 2: Operativo (1/2)}
  \footnotesize
  \begin{block}{Problema 5.3.6 --- Prueba de Hipótesis para $\sigma^2$ (Enunciado)}
    Una máquina embotelladora debe tener varianza $\sigma_0^2=10$. Una muestra de $n=16$ botellas produce $S^2=18.2$. Calcule $(n-1)S^2/\sigma_0^2$ y decida si se rechaza $H_0:\sigma^2=10$ a nivel $\alpha=0.05$, usando $\chi^2_{15,0.95}\approx 24.996$.
  \end{block}

  \vspace{0.15cm}
  \pause
  \begin{alertblock}{Estrategia}
    \scriptsize
    Compare el estadístico calculado contra el valor crítico dado.
  \end{alertblock}
\end{frame}

% Slide 13B: Ejercicio Nivel 2 (Resolución)
\begin{frame}{Ejercicio en Clase --- Nivel 2: Operativo (2/2)}
  \scriptsize
  Calculamos el estadístico y comparamos: \pause

  \vspace{0.05cm}
  \begin{block}{Cálculo}
    $$\frac{(n-1)S^2}{\sigma_0^2} = \frac{15\cdot 18.2}{10} = 27.3$$
  \end{block}

  \vspace{0.05cm}
  \pause
  \begin{alertblock}{Decisión}
    \scriptsize
    Como $27.3 > 24.996 = \chi^2_{15,0.95}$, \textbf{rechazamos} $H_0:\sigma^2=10$ a nivel $\alpha=0.05$: hay evidencia estadística de que la varianza real de la máquina excede el objetivo de $10$.
  \end{alertblock}
\end{frame}

% Slide 14A: Ejercicio Nivel 3 (Enunciado)
\begin{frame}{Ejercicio en Clase --- Nivel 3: Analítico (1/2)}
  \footnotesize
  \begin{block}{Problema 5.3.7 --- Derivación de $E(\chi^2_\nu)$ y $\Var(\chi^2_\nu)$ (Enunciado)}
    Partiendo de $\chi^2_\nu=\text{Gamma}(\nu/2,2)$ y de $\mu_X=\alpha\beta$, $\sigma_X^2=\alpha\beta^2$ (Sección 04.07), derive algebraicamente que $E(\chi^2_\nu)=\nu$ y $\Var(\chi^2_\nu)=2\nu$.
  \end{block}

  \vspace{0.1cm}
  \pause
  \begin{alertblock}{Estrategia}
    \scriptsize
    Sustituya $\alpha=\nu/2$ y $\beta=2$ directamente en las fórmulas generales de la Gamma.
  \end{alertblock}
\end{frame}

% Slide 14B: Ejercicio Nivel 3 (Resolución)
\begin{frame}{Ejercicio en Clase --- Nivel 3: Analítico (2/2)}
  \scriptsize
  Sustituimos $\alpha=\nu/2$, $\beta=2$: \pause

  \vspace{0.05cm}
  \begin{block}{Derivación}
    \begin{align*}
      E(\chi^2_\nu) = \alpha\beta = \frac{\nu}{2}\cdot 2 = \nu, \qquad
      \Var(\chi^2_\nu) = \alpha\beta^2 = \frac{\nu}{2}\cdot 4 = 2\nu. \quad \qedhere
    \end{align*}
  \end{block}

  \vspace{0.05cm}
  \pause
  \begin{alertblock}{Interpretación}
    \scriptsize
    Esta derivación conecta directamente con la Sección 04.07: toda propiedad ya demostrada para la familia Gamma general se hereda automáticamente por sus casos particulares, incluyendo la chi-cuadrada.
  \end{alertblock}
\end{frame}

% Slide 15A: Ejercicio Nivel 4 (Enunciado)
\begin{frame}{Ejercicio en Clase --- Nivel 4: Desafiante (1/2)}
  \footnotesize
  \begin{block}{Problema 5.3.9 --- De la Normal a la $t$ de Student (Enunciado)}
    Usando el Teorema de Fisher, muestre que $T = \dfrac{\bar X-\mu}{S/\sqrt n} = \dfrac{Z}{\sqrt{\chi^2_{n-1}/(n-1)}}$, donde $Z=(\bar X-\mu)/(\sigma/\sqrt n)\sim N(0,1)$. ¿Por qué es relevante esta descomposición cuando $\sigma$ es desconocida?
  \end{block}

  \vspace{0.15cm}
  \pause
  \begin{alertblock}{Estrategia}
    \scriptsize
    Divida numerador y denominador de $T$ entre $\sigma/\sqrt n$, y reconozca $S^2/\sigma^2$ en términos del Teorema de Fisher.
  \end{alertblock}
\end{frame}

% Slide 15B: Ejercicio Nivel 4 (Resolución)
\begin{frame}{Ejercicio en Clase --- Nivel 4: Desafiante (2/2)}
  \scriptsize
  Dividimos numerador y denominador: \pause

  \vspace{0.05cm}
  \begin{block}{Derivación}
    \scriptsize
    \begin{align*}
      T = \frac{\bar X-\mu}{S/\sqrt n} = \frac{(\bar X-\mu)/(\sigma/\sqrt n)}{S/\sigma} = \frac{Z}{\sqrt{S^2/\sigma^2}} = \frac{Z}{\sqrt{\chi^2_{n-1}/(n-1)}}. \quad \qedhere
    \end{align*}
  \end{block}

  \vspace{0.05cm}
  \pause
  \begin{alertblock}{Relevancia}
    \scriptsize
    Cuando $\sigma$ es desconocida, $T$ ya no es Normal exacta: es el cociente de una Normal estándar entre la raíz de una chi-cuadrada independiente dividida entre sus grados de libertad --- precisamente la definición de la distribución $t$ de Student que estudiaremos en la Sección 05.04.
  \end{alertblock}
\end{frame}

% Slide 16: Cierre
\begin{frame}{Cierre de Sección: Chi-Cuadrada y Teorema de Fisher}
  \begin{reflexion}
    El Teorema de Fisher es el puente entre la teoría de la distribución Normal y la inferencia práctica con $\sigma$ desconocida: nos dice exactamente cómo se distribuye $S^2$, y que lo hace de forma independiente de $\bar X$. Sin este resultado, no podríamos construir intervalos exactos para $\sigma^2$ ni derivar la distribución $t$ de Student.
  \end{reflexion}

  \vspace{0.2cm}
  \pause
  \begin{block}{Lecciones Clave}
    \begin{itemize}\itemsep=0.08cm
      \item \textbf{Chi-cuadrada:} $\chi^2_\nu=\sum Z_i^2$, con $E=\nu$, $\Var=2\nu$, caso particular de la Gamma.
      \item \textbf{Teorema de Fisher:} para muestras normales, $\bar X \perp S^2$ y $(n-1)S^2/\sigma^2\sim\chi^2_{n-1}$.
      \item \textbf{Aplicación:} intervalos de confianza exactos (asimétricos) y pruebas de hipótesis para $\sigma^2$.
    \end{itemize}
  \end{block}
\end{frame}

% Slide 17: Perspectiva Modular
\begin{frame}{Perspectiva Modular --- El Próximo Reto}
  \scriptsize
  \begin{retoclase}
    Hoy vimos que $T=(\bar X-\mu)/(S/\sqrt n)$ se descompone como el cociente de una Normal estándar entre la raíz de una chi-cuadrada dividida entre sus grados de libertad. La Sección 05.04 formalizará esta construcción como la distribución $t$ de Student, la herramienta exacta para inferencia sobre $\mu$ cuando $\sigma$ es desconocida y la muestra es pequeña.
  \end{retoclase}

  \vspace{0.08cm}
  \pause
  \begin{alertblock}{Preparación Recomendada}
    \begin{itemize}\itemsep=0.06cm
      \item Repasa la demostración del Problema 5.3.9: es exactamente la construcción de la $t$ de Student.
      \item Reflexiona sobre por qué la $t$ tiene colas más pesadas que la Normal para $n$ pequeño.
    \end{itemize}
  \end{alertblock}
\end{frame}

\end{document}
